\chapter{Conclusion and Future Scope}
\label{chap:conclusion}

\section{Summary}
This academic investigation has systematically explored the analytical and experimental dimensions of computational complexity, grounded in Kenneth H. Rosen's \emph{Discrete Mathematics and Its Applications} (8th Edition), Chapter 3 \cite{rosen2019discrete}. Through rigorous mathematical modeling, deterministic telemetry instrumentation, and extensive statistical benchmarking, we examined Linear Search, Binary Search, Bubble Sort, and Insertion Sort across a wide spectrum of input scales and topologies.

\section{Major Findings}
\begin{enumerate}
    \item \textbf{Mathematical Model Exactness}: Closed-form operation-count formulas derived from discrete mathematics ($T(n) = n$, $T(n) = 2(\lfloor \log_2 n \rfloor + 1)$, $T(n) = n(n-1)/2$, $T(n) = n - 1$) accurately model empirical counter telemetry with \textbf{0.00\% discrepancy}.
    \item \textbf{Asymptotic Scaling Domination}: At $n = 100\,000$, Binary Search demonstrates an empirical speedup of approximately \textbf{$1\,946\times$} over Linear Search for unsuccessful queries, confirming that logarithmic interval halving decisively eclipses linear scanning.
    \item \textbf{Inversion Equivalence Verification}: We verified empirically that Bubble Sort adjacent swaps and Insertion Sort right shifts are identical ($2\,296\,691$ operations at $n = 3\,000$) on the same random permutation, confirming that both algorithms resolve the identical number of permutation inversions $I(A)$.
    \item \textbf{Adaptive Linear Sorting}: Insertion Sort operates in linear time $\Theta(n)$ ($0.32\text{ ms}$ at $n = 3\,000$) on pre-sorted data, outperforming fixed-pass Bubble Sort by a factor of \textbf{$769\times$}.
\end{enumerate}

\section{Achievement of Objectives}
All milestone objectives established in Milestone M1 have been successfully realized:
\begin{itemize}
    \item Full mathematical derivations aligned with Rosen §3.1, §3.2, and §3.3.
    \item Dual-mode algorithm implementations with unit-tested correctness ($64$ automated pytest tests).
    \item Telemetry-isolated empirical benchmarking with $k = 30$ repetitions.
    \item Generation of nine publication-grade figures in dual vector PDF and 300 DPI raster PNG formats.
    \item A fully executed, self-contained interactive Jupyter Research Notebook.
    \item A comprehensive, modular academic monograph typeset in LaTeX.
\end{itemize}

\section{Conclusion}
The study demonstrates that discrete mathematical principles provide an extraordinarily reliable framework for predicting real-world algorithm behavior. As computational datasets continue to expand, formal algorithm analysis remains the most vital tool for building scalable, high-performance software systems.

\section{Future Scope}
Future research directions extending this study include:
\begin{enumerate}
    \item \textbf{Higher-Order Complexity Classes}: Extending empirical benchmarking to $\mathcal{O}(n \log n)$ algorithms (Merge Sort, Quick Sort, Heap Sort) and linear non-comparison sorts (Radix Sort, Counting Sort).
    \item \textbf{Cache-Oblivious Algorithms}: Investigating how memory hierarchy caching, TLB misses, and vector SIMD instructions influence constant factors at scales exceeding $n = 10^7$.
    \item \textbf{Parallel & Distributed Implementations}: Evaluating parallel bisection and multi-threaded sorting networks across multi-core processor architectures.
\end{enumerate}
